% ============================================================
%  Paper 10: Conformal Higgs Mass from Golden Ratio and Lattice Capacity
%  Target: RevTeX 4-2 Compilation (SDGFT Series)
% ============================================================
\documentclass[amsmath,amssymb,aps,prd,twocolumn,superscriptaddress,nofootinbib]{revtex4-2}

\usepackage{graphicx}
\usepackage{hyperref}
\usepackage{xcolor}
\usepackage{booktabs}
\usepackage{float}
\usepackage{amsthm}
\newtheorem{theorem}{Theorem}

% ── SDGFT shared macros ──────────────────────────────────────
\usepackage{sdgft-macros}

% ── Document ─────────────────────────────────────────────────
\begin{document}

\preprint{SDGFT-2026-10}

\title{Conformal Higgs Mass from Golden Ratio and Lattice Capacity}

\author{David A. Besemer}
\email{david.besemer@sdgft.org}
\affiliation{Independent Researcher}

\date{\today}

\begin{abstract}
We calculate the mass of the Higgs boson from the conformal geometry of the F4 lattice. The physical mass is derived from the tree-level Z-boson mass, corrected by the topological conformal scale factor (1 + Delta/phi). This yields a Higgs mass of MH ~ 125.11 GeV, in agreement with the current experimental average.
\end{abstract}

\maketitle

\section{Introduction}
The discovery of the Higgs boson at ~125 GeV completed the Standard Model. We derive its mass from the conformal ratio of the 24-cell lattice.

\section{Theoretical Formulation}
Here we formulate the core mathematical relations governing the system:
\begin{equation}
  M_H = M_Z \sqrt{\frac{23}{12} \frac{37}{48}} \left(1 + \frac{\DD}{\gp}\right)
\end{equation}
\begin{equation}
  M_H \approx 91.1876 \times 1.21544 \times 1.12876 \approx 125.11 \text{ GeV}
\end{equation}
\section{Conformal Correction Factor}
The correction factor (1 + Delta/phi) represents the scale factor between the F4 lattice tension and the golden ratio, which governs electroweak symmetry breaking.

\section{Conclusion}
The Higgs mass is not a free parameter but is fixed by the conformal structure of the electroweak sector.



\begin{thebibliography}{99}
\bibitem{Goldstein_Paper01} D. A. Besemer, ``Axiomatic Foundations of SDGFT,'' SDGFT-2026-01 (in preparation).
\bibitem{Goldstein_Paper04} D. A. Besemer, ``Effective Spectral Dimension and the Banach Fixed-Point Theorem in SDGFT,'' SDGFT-2026-04.
\bibitem{Goldstein_Code} D. A. Besemer, \texttt{sdgft} Python package, \url{https://github.com/david-besemer/sdgft} (2026).
\end{thebibliography}

\end{document}
